\section{Same fields, odd lengths, and the physical gap}\label{sec:conclusion}
The estimates obtained so far concern even moments of the opposite-field
chain. We now use the rotation-invariant leading line to control the
same-field partition function at every odd length in a dyadic interval.

\begin{lemma}[Signed boundary comparison]\label{co:signed}
For $j\ge0$, an even $M\ge2$, and any integer $n\ge M$, define
$Q_j(n)=Y_n^{\rm s}(b_j)/\lambda_j^n$. Then
\begin{equation}\label{co:error}
 |Q_j(n)-c_j|\le B_M(j)-c_j.
\end{equation}
\end{lemma}
\begin{proof}
The leading spatial line is fixed by $U_P$, and its eigenvalue is
$+\lambda_j$. Thus its contribution to $Q_j(n)$ is exactly $c_j$.
For every other eigenvalue $\xi$ of $X_{b_j}$, let $P_\xi$ be its
spectral projection. The unitary $U_P$ commutes with that projection,
so Cauchy--Schwarz gives
\[
 |\langle P_\xi v_{b_j},U_P P_\xi v_{b_j}\rangle|
 \le\norm{P_\xi v_{b_j}}^2.
\]
Consequently the absolute residual is at most
\[
 \sum_{\xi\ne\lambda_j}
       (|\xi|/\lambda_j)^n\norm{P_\xi v_{b_j}}^2
 \le B_M(j)-c_j.
\]
The series converges absolutely, being bounded by $\norm{v_{b_j}}^2$.
The kernel of $X_{b_j}$ contributes zero because all exponents are
positive. Negative spatial eigenvalues and odd powers are therefore
included in the estimate.
\end{proof}

One concentrated even moment therefore controls the normalized same-field
pairing $Q_j(n)$ at every longer length, including odd lengths. To compare
two temperatures, we use the common even reference length $M_j=k_{j+1}$.
The lower bound for $r_j^{M_j}$ will cover the interval $M_j\le n\le2M_j$.

\begin{proposition}\label{co:purity}
For $j\ge5$ and every integer $n$ with $M_j\le n\le2M_j$,
\begin{equation}\label{co:purity-bound}
 \frac{Y_n^{\rm s}(b_{j+1})}{Y_n^{\rm s}(b_j)^2}\ge1-4x_j.
\end{equation}
\end{proposition}
\begin{proof}
Put $M=M_j$. At temperature $b_j$, Lemma~\ref{co:signed} gives
$Q_j(n)\le B_M(j)$. This upper bound can be squared because the
physical partition function, and hence $Q_j(n)$, is positive.
At temperature $b_{j+1}$, use the same even length $M=k_{j+1}$:
\[
 Q_{j+1}(n)\ge2c_{j+1}-B_M(j+1)
 \ge(1-2x_{j+1})B_M(j+1)>0.
\]
The last step uses Proposition~\ref{in:quadratic}. It follows that
\begin{align}
 \frac{Y_n^{\rm s}(b_{j+1})}{Y_n^{\rm s}(b_j)^2}
 &\ge(1-2x_{j+1})A(M,j)r_j^{n-M} \notag\\
 &\ge(1-2x_{j+1})(1-x_j)(1-u_j)(1-u_{j+1}).
                                                    \label{co:product}
\end{align}
Indeed $0<r_j\le1$ and $0\le n-M\le M$ give
$r_j^{n-M}\ge r_j^M\ge(1-u_j)(1-u_{j+1})$.
Using $x_{j+1}=10x_j^2$, $u_j\le x_j$, and
$u_{j+1}\le6x_j^2$, the defect of the last product is at most
\[
 20x_j^2+x_j+x_j+6x_j^2
 =2x_j+26x_j^2\le4x_j,
\]
as $x_j\le.01$. This proves \eqref{co:purity-bound}.
\end{proof}

\begin{proof}[Proof of Theorem~\ref{thm:main}]
Fix an odd $n$ in $[M_j,2M_j]$ with $j\ge5$.
Let $p_0\ge p_1\ge\cdots$ be the Gibbs probabilities of $H_n^{\rm s}$
at inverse temperature $b_j$, listing energies with multiplicity.
The left side of \eqref{co:purity-bound} is $\sum_i p_i^2$.
Since $\sum_i p_i^2\le p_0\sum_i p_i=p_0$, we have
$p_0\ge1-4x_j>1/2$. In particular the physical ground state is
simple, and
\begin{equation}\label{co:ratio}
 e^{-b_j(E_1-E_0)}=\frac{p_1}{p_0}
 \le\frac{1-p_0}{p_0}
 \le\frac{4x_j}{1-4x_j}
 <10^{-2^{j-5}}.
\end{equation}
For the strict last inequality, use \eqref{in:xexplicit} and
$4/[10(1-.04)]<1$. Because $b_j=392\,2^{j-5}$,
taking logarithms yields $E_1-E_0>\log(10)/392$.

Finally $M_5=1920$ and $M_{j+1}=2M_j$. The intervals
$[M_j,2M_j]$, $j\ge5$, cover every integer $n\ge1920$.
Every odd length $n=2L+1$ with $L\ge960$ is therefore covered.
Relabeling the sites gives precisely \eqref{eq:target}, with the
single field $h=3/5$ fixed in \eqref{tr:constants}.
\end{proof}
